\documentclass[11pt]{article}

\usepackage[margin=1.08in]{geometry}
\usepackage[T1]{fontenc}
\usepackage{lmodern}
\IfFileExists{microtype.sty}{\usepackage{microtype}}{}
\usepackage{amsmath,amssymb,amsthm,mathtools}
\usepackage{booktabs}
\usepackage{array}
\usepackage{xcolor}
\usepackage{hyperref}
\hypersetup{
  pdfauthor={Joshua Mark},
  pdfsubject={Sublogarithmic Hadamard existence, sparse representations and formal verification},
  pdfkeywords={Hadamard matrices, Golay complementary pairs, multi-base representations, Lean},
  colorlinks=true,
  linkcolor=blue!55!black,
  citecolor=blue!55!black,
  urlcolor=blue!55!black,
  pdftitle={Sparse Multi-Base Representations and Asymptotic Existence of Hadamard Matrices}
}

\newtheorem{theorem}{Theorem}[section]
\newtheorem{proposition}[theorem]{Proposition}
\newtheorem{corollary}[theorem]{Corollary}
\newtheorem{lemma}[theorem]{Lemma}
\newtheorem{remark}[theorem]{Remark}
\newtheorem{definition}[theorem]{Definition}

\newcommand{\Hcal}{\mathcal{H}}
\newcommand{\N}{\mathbb{N}}
\newcommand{\R}{\mathbb{R}}
\newcommand{\eps}{\varepsilon}
\newcommand{\wt}{\operatorname{wt}}
\newcommand{\ee}{\mathrm{e}}
\newcommand{\scale}{\Phi}
\setlength{\emergencystretch}{2em}

\title{\textbf{Sparse Multi-Base Representations and\\Asymptotic Existence of Hadamard Matrices}}
\author{Joshua Mark\thanks{ORCID: \href{https://orcid.org/0009-0003-5905-2382}{0009-0003-5905-2382}. AI tools assisted with proof development, formalization, exposition, and manuscript preparation. The author is responsible for the submitted work.}}
\date{4 November 2026}

\begin{document}
\maketitle

\begin{abstract}
For an odd positive integer $p$, let
\[
\tau(p):=\min\{t\in\N_0: \text{a Hadamard matrix of order }2^t p\text{ exists}\}.
\]
We combine the sparse-representation theorem of Krenn, Suppakitpaisarn and Wagner with Craigen's complementary-sequence construction for Hadamard matrices.  The resulting transfer gives the uniform bound
\[
\tau(p)=O\!\left(\frac{\log p}{\log\log p}\right),
\]
uniformly over odd $p$ as $p$ tends to infinity. In particular, the constructed exponent is sublogarithmic.  The key observation is that every integer $n>2$ is a sum of $O(\!\log n/\log\log n)$ distinct numbers $2^a10^b$, while every $2^a10^b$ is the length of a binary Golay complementary pair.

We also address an additive-number-theory problem posed by Livinskyi.  If $S(b_1,\ldots,b_r)$ is the multiplicative monoid generated by fixed multiplicatively independent integers $b_1=2,b_2,\ldots,b_r$, and $\ell_m(S)$ is the smallest positive integer not expressible as a sum of at most $m$ elements of $S$, then
\[
\exp(c m\log m)\le \ell_m(S)\le
\exp\bigl((r-1+o(1))m\log m\bigr)
\]
for some $c>0$.  In particular,
$\ell_m(S(2,10,26))=\exp(\Theta(m\log m))$.
The same construction gives a uniform bound on distances to larger Hadamard orders, and hence the consecutive-order gap estimate
\[
\gamma(x)\le \exp\!\left(O\!\left(\frac{\log x}{\log\log x}\right)\right)=x^{o(1)}
\]
for consecutive Hadamard orders.  Combining this with the random-rounding argument of Alexeev, Jasper and Mixon yields, for every $\delta>0$, the asymptotic approximate-Hadamard bound
\[
\kappa(n)\le 1+n^{-1/4+\delta}
\]
for all sufficiently large $n$. The accompanying Lean development proves the main bound and its stated consequences without additional arithmetic hypotheses.
\end{abstract}

\section{Introduction}
A Hadamard matrix of order $n$ is a matrix $H\in\{\pm1\}^{n\times n}$ satisfying
\[
HH^{\mathsf T}=nI_n.
\]
Except for orders $1$ and $2$, divisibility by $4$ is necessary.  The Hadamard conjecture asserts that this necessary condition is sufficient: a Hadamard matrix should exist at every order $4k$.

A classical asymptotic relaxation fixes an odd integer $p$ and asks how many powers of two must be adjoined before a Hadamard order is guaranteed.  Wallis proved the first general logarithmic bound in 1976~\cite{Wallis1976}; Craigen improved the coefficient using signed groups and complementary sequences~\cite{Craigen1995}.  More recently, Du and Jiang proved that for every odd $m$ there is a Hadamard matrix of order $2^t m$ with
\begin{equation}\label{eq:dujiang}
 t=6\left\lfloor \frac{1}{40}\log_2 m\right\rfloor+10
   =\frac{3}{20}\log_2 m+O(1).
\end{equation}
Their paper also points out that one of the stronger numerical bounds stated in Livinskyi's 2012 thesis relies on an earlier theorem that is incorrect~\cite{DuJiang2024}.  The argument developed here does not use that result: it uses only the binary-Golay transfer theorem originating with Craigen and quoted as Theorem~3.10 in Livinskyi's thesis~\cite{Craigen1995,Livinskyi2012}.

Livinskyi also proposed a different route.  Write
\[
S(n_1,\ldots,n_k)=\{n_1^{a_1}\cdots n_k^{a_k}:a_i\in\N_0\}
\]
and let $\ell_m(S)$ be the smallest positive integer not representable as a sum of at most $m$ elements of $S$.  Since
\[
S(2,10,26)
\]
consists of binary Golay lengths, Livinskyi observed that strong lower bounds on $\ell_m(S(2,10,26))$ would translate into stronger bounds for $\tau(p)$, but no asymptotic estimate for this covering function was obtained there~\cite{Livinskyi2012}.

The missing ingredient appears in a different part of number theory.  Krenn, Suppakitpaisarn and Wagner proved that, for any fixed set of at least two multiplicatively independent bases satisfying a mild representability assumption, the worst-case minimal Hamming weight of a multi-base representation is
\[
\Theta\!\left(\frac{\log n}{\log\log n}\right).
\]
Their upper bound follows by analyzing a greedy construction~\cite{Krenn2020}. In our application the digit set is $\{0,1\}$ from the outset; binary expansion verifies the required representability assumption.  Applying their theorem to the bases $2$ and $10$ produces sparse sums of binary Golay lengths, exactly the input required by Craigen's Hadamard construction.

The results established here are:
\begin{enumerate}
\item a sparse Golay decomposition of every positive integer using $O(\log n/\log\log n)$ binary Golay lengths;
\item the uniform Hadamard-existence estimate $\tau(p)=O(\log p/\log\log p)$ for all odd $p$;
\item the order of growth $\log \ell_m(S(2,10,26))=\Theta(m\log m)$, answering Livinskyi's proposed additive-number-theory direction at the level of order of magnitude;
\item a quantitative subpolynomial bound on gaps between Hadamard orders; and
\item the unconditional limiting exponent of the gap-based method of Alexeev--Jasper--Mixon~\cite{Alexeev2025}, for approximate Hadamard matrices: every $0<\alpha<1/4$ is available, compared with the range $0<\alpha<17/92$ in their stated bound.
\end{enumerate}

The Hadamard-existence result is a deduction from established construction
and sparse-representation theorems. We make no claim here to priority for
the qualitative subpolynomial gap statement: Brent, Osborn and Smith
already attributed such a statement to unpublished work of Smith
\cite{Brent2013}. The quantitative rate and the covering estimates are stated
separately from that historical observation.

\paragraph{Notation and quantifiers.}
Write $\N_0=\{0,1,2,\ldots\}$ and $\N=\{1,2,\ldots\}$. All logarithms
without a subscript are natural. For real $x>\ee$, put
\[
\scale(x)=\frac{\log x}{\log\log x}.
\]
This function is positive on that domain and increases for $x>\ee^{\ee}$.
The expression $f(x)=O(g(x))$ means that one positive constant and one
threshold bound $|f(x)|$ by that constant times $|g(x)|$ at every later
input. A bound described as uniform over odd inputs uses the same constant
and threshold for all such inputs. The notation $o(1)$ denotes a quantity
tending to zero. Constants may depend on a fixed collection of bases or a
fixed parameter such as $\varepsilon$, but not on the input tending to infinity.

\section{Golay pairs, complementary sequences, and multi-base representations}

\subsection{Binary Golay lengths}
For a binary sequence $A=(a_0,\ldots,a_{n-1})\in\{\pm1\}^n$, write its aperiodic autocorrelation at positive shift $s$ as
\[
N_A(s)=\sum_{j=0}^{n-s-1}a_j a_{j+s}.
\]
Set $N_A(s)=0$ for $s\ge n$, in accordance with the empty-sum convention. A pair $(A,B)$ of binary sequences of the same length $n$ is a \emph{Golay complementary pair} if
\[
N_A(s)+N_B(s)=0\qquad(1\le s<n).
\]
The corresponding $n$ is called a binary Golay length.  Golay introduced complementary series~\cite{Golay1961}; Turyn's multiplication construction implies the following standard family~\cite{Turyn1974}.

\begin{proposition}[Golay--Turyn]\label{prop:golaylengths}
For every $a,b,c\in\N_0$, there exists a binary Golay complementary pair of length
\[
2^a10^b26^c.
\]
The length $1$ is understood as the trivial pair.
\end{proposition}

Only the subfamily $2^a10^b$ will be needed for the main theorem.  The larger monoid $S(2,10,26)$ is useful for the covering problem in Section~\ref{sec:covering}.

\subsection{Craigen's transfer to Hadamard matrices}
The following construction is due to Craigen~\cite{Craigen1995}; the convenient formulation below is Theorem~3.10 of Livinskyi~\cite{Livinskyi2012}.

\begin{theorem}[Craigen]\label{thm:craigen-transfer}
Let $r\ge1$. Suppose that
\[
A_1,B_1,\ldots,A_r,B_r
\]
are $2r$ binary sequences whose total aperiodic autocorrelation is zero at every positive shift, with $A_i$ and $B_i$ both of length $l_i$.  Put $L=\sum_{i=1}^r l_i$.  Then there exists a Hadamard matrix of order
\[
4^{r+1}(2L+1).
\]
\end{theorem}

If each $(A_i,B_i)$ is a Golay pair, the hypothesis is automatic because the autocorrelation sums vanish pairwise.

\subsection{Sparse multi-base representations}
Fix integers $b_1,\ldots,b_s>1$.  Their power-product monoid is
\[
\mathcal B=\{b_1^{a_1}\cdots b_s^{a_s}:a_i\in\N_0\}.
\]
For a finite digit set $D\subset\N_0$, a multi-base representation of $n$ is
\[
 n=\sum_{B\in\mathcal B} d_BB,
 \qquad d_B\in D,
\]
with only finitely many nonzero coefficients.  Its Hamming weight is the number of nonzero $d_B$.  The bases are \emph{multiplicatively independent} if
\[
b_1^{z_1}\cdots b_s^{z_s}=1,\qquad z_i\in\mathbb Z,
\]
implies $z_1=\cdots=z_s=0$.

We use the following result of Krenn, Suppakitpaisarn and Wagner~\cite{Krenn2020}.

\begin{theorem}[Krenn--Suppakitpaisarn--Wagner]\label{thm:ksw}
Let $s\ge2$.  Suppose $b_1,\ldots,b_s$ are multiplicatively independent and the finite nonnegative digit set $D$, containing $0$, permits a multi-base representation of every positive integer.  There are constants $K_1,K_2>0$ such that:
\begin{enumerate}
\item every $n>2$ has a representation of Hamming weight at most
\[
K_1\frac{\log n}{\log\log n};
\]
\item for infinitely many $n$, every such representation has Hamming weight at least
\[
K_2\frac{\log n}{\log\log n}.
\]
\end{enumerate}
In the applications below we take $D=\{0,1\}$ and verify its representability hypothesis separately; no conversion of an arbitrary digit set into $\{0,1\}$ is asserted.
\end{theorem}

The constants depend on the bases and digits.  The proof is asymptotic; in particular, the preprocessing threshold in the greedy algorithm may be very large~\cite{Krenn2020}.

\section{Sparse decompositions into binary Golay lengths}
For the Hadamard construction, only the following specialization of Theorem~\ref{thm:ksw} is required.

\begin{lemma}\label{lem:independence}
The integers $2$ and $10$ are multiplicatively independent.
\end{lemma}
\begin{proof}
If $2^a10^b=1$ for $a,b\in\mathbb Z$, then
\[
2^{a+b}5^b=1.
\]
Unique factorization forces $b=0$, and then $a=0$.
\end{proof}

\begin{theorem}[Sparse Golay decomposition]\label{thm:sparse-golay}
There is a constant $K>0$ such that every integer $q>2$ can be written
\begin{equation}\label{eq:sparse-golay}
q=\sum_{j=1}^{r}2^{a_j}10^{b_j},
\qquad a_j,b_j\in\N_0,
\end{equation}
with the summands distinct and
\[
r\le K\frac{\log q}{\log\log q}.
\]
Every summand in \eqref{eq:sparse-golay} is a binary Golay length.
\end{theorem}

\begin{proof}
Apply Theorem~\ref{thm:ksw} with bases $(2,10)$ and digit set $D=\{0,1\}$.  Multiplicative independence follows from Lemma~\ref{lem:independence}.  The representability hypothesis holds because ordinary binary expansion is already a multi-base representation using only the power-products $2^a10^0$.  The theorem therefore gives a representation with $O(\log q/\log\log q)$ nonzero $0/1$ digits, which is exactly \eqref{eq:sparse-golay}.  Proposition~\ref{prop:golaylengths} gives a binary Golay pair at each such length.
\end{proof}

\begin{remark}\label{rem:morebases}
The same argument works with any fixed collection of multiplicatively independent binary Golay generators containing $2$.  The pair $(2,10)$ is enough for the asymptotic order proved here and keeps the transfer transparent.
\end{remark}

\section{A sublogarithmic exponent for asymptotic Hadamard existence}
Let $\Hcal\subset\N$ be the set of Hadamard orders. We first check that
the minimum used in the main statement exists.

\begin{lemma}[Existence for every odd multiplier]\label{lem:existence}
For every positive odd integer $p$, some $t\in\N_0$ satisfies $2^tp\in\Hcal$.
\end{lemma}
\begin{proof}
For $p=1$, use the order-one matrix. If $p>1$, put $q=(p-1)/2\ge1$.
Binary expansion expresses $q$ as a sum of distinct powers of two.
Choose a Golay pair at each of those lengths using
Proposition~\ref{prop:golaylengths}, and apply
Theorem~\ref{thm:craigen-transfer}. If the sum has $r$ terms, the resulting
order is $4^{r+1}p=2^{2r+2}p$.
\end{proof}

Consequently the nonempty set of admissible exponents has a least element:
\[
\tau(p)=\min\{t\in\N_0:2^tp\in\Hcal\}.
\]
Here ``admissible'' means that the indicated matrix exists; it does not mean
that the exponent is the smallest possible.

\begin{theorem}[Main Hadamard bound]\label{thm:main}
There exist a real constant $C>0$ and an integer $N\ge3$ such that, for every
odd integer $p\ge N$, there are $t\in\N_0$ and a Hadamard matrix of order
$2^tp$ satisfying
\begin{equation}\label{eq:main}
\tau(p)\le t\le C\frac{\log p}{\log\log p}.
\end{equation}
Thus $\tau(p)=O(\log p/\log\log p)$ over odd inputs. The constructed order
has the form
\[
2^tp\le p\exp\!\left(\frac{(C\log2)\log p}{\log\log p}\right)
=p^{1+(C\log2)/\log\log p}=p^{1+o(1)}.
\]
Neither primality of $p$ nor a Hadamard conjecture is assumed.
\end{theorem}

\begin{proof}
For odd $p\ge7$, put $q=(p-1)/2>2$. Theorem~\ref{thm:sparse-golay}
gives distinct Golay lengths $l_1,\ldots,l_r$ such that
\[
q=l_1+\cdots+l_r,\qquad r\le K\scale(q).
\]
Choose a Golay pair of length $l_i$ for each $i$. The total autocorrelation
of these pairs vanishes at every positive shift, including shifts beyond
an individual length (where its contribution is an empty sum).
Theorem~\ref{thm:craigen-transfer} therefore constructs order
\[
4^{r+1}(2q+1)=2^{2r+2}p.
\]
Set $t=2r+2$. Since $\scale((p-1)/2)/\scale(p)\to1$, eventually
$\scale(q)\le2\scale(p)$; also $\scale(p)\ge1$ eventually. Hence
\[
t\le2K\scale(q)+2\le(4K+2)\scale(p).
\]
Choose $C=4K+2$ and enlarge the threshold to cover the preceding conditions.
The minimum $\tau(p)$ is at most this constructed exponent.
\end{proof}

\begin{remark}[An exact finite illustration]
Take $p=47$, so $q=23=20+2+1$. These three distinct lengths belong to
$S(2,10)$, giving $r=3$ and $t=8$. The transfer constructs order
$2^8\cdot47=12032$. This is an upper bound on $\tau(47)$, not a computation
of its minimum or a claim that $47$ is beyond the asymptotic threshold.
\end{remark}

\begin{remark}[Comparison with the recent logarithmic bound]
Du and Jiang's formula \eqref{eq:dujiang} is explicit and is likely far stronger at computationally accessible sizes.  Theorem~\ref{thm:main} is instead an asymptotic statement: since
\[
\frac{\log p}{\log\log p}=o(\log p),
\]
it eventually improves every positive constant multiple of $\log p$.  No useful crossover point follows from the present proof because the constants inherited from Theorem~\ref{thm:ksw} are not explicit in a practical form.
\end{remark}

\begin{remark}[A limitation of the fixed multi-base mechanism]\label{rem:optimality}
The lower-bound half of Theorem~\ref{thm:ksw} shows that, for the fixed bases $(2,10)$ and digits $\{0,1\}$, infinitely many integers require
$\Omega(\log n/\log\log n)$ distinct power-products.  Allowing repeated summands does not circumvent this order: whenever two equal elements of $S(2,10)$ occur, they may be merged into their double, which remains in $S(2,10)$; iterating produces a representation by distinct elements without increasing the number of terms.  Thus the pair count in Theorem~\ref{thm:sparse-golay} has the correct worst-case order within this fixed-monoid representation framework. This is not a lower bound of that order on $\tau(p)$: other matrix constructions may use smaller exponents.  A further asymptotic improvement through this representation-and-transfer mechanism would require additional structure---for example, a richer family of Golay lengths, a different additive representation theorem, or a Hadamard transfer whose power-of-two overhead grows more slowly than linearly in the number of complementary pairs.
\end{remark}

\section{Livinskyi's covering function}\label{sec:covering}
For $m\in\N_0$ and $S\subset\N$, define the covering function of Livinskyi, whenever the set below is nonempty, by
\begin{equation}\label{eq:ell-def}
\ell_m(S):=\min\{n\ge1:n\text{ is not a sum of at most }m\text{ elements of }S\}.
\end{equation}
Repetitions are allowed in this definition, and the empty sum represents zero. He singled out $S(2,10,26)$ because it is contained in the binary Golay lengths~\cite{Livinskyi2012}. For the fixed finitely generated monoids below, a missing positive integer exists for every fixed $m$: the counting argument in the proof bounds the number of available sums up to $N$ by a power of $\log N$, which is smaller than $N$ for large $N$.  The multi-base theorem determines the order of growth of this function up to constants in the exponent.

\begin{theorem}[Covering by a fixed multiplicative monoid]\label{thm:ell-general}
Let $r\ge2$, let $b_1=2,b_2,\ldots,b_r>1$ be multiplicatively independent, and put
\[
S=S(b_1,\ldots,b_r)=\{b_1^{a_1}\cdots b_r^{a_r}:a_i\in\N_0\}.
\]
Then there exists $c>0$ such that, for every $\eta>0$, there is an integer $M$ for which every integer $m\ge M$ satisfies
\begin{equation}\label{eq:ell-bounds}
\exp(c m\log m)
\le \ell_m(S)
\le \exp\bigl((r-1+\eta)m\log m\bigr).
\end{equation}
In particular,
\[
\log\ell_m(S)=\Theta(m\log m).
\]
\end{theorem}

\begin{proof}
\emph{Lower bound.}
Apply the upper-bound half of Theorem~\ref{thm:ksw} with digit set $\{0,1\}$.  The representability assumption holds because $2$ is one of the bases.  Hence every $n>2$ is a sum of at most
\[
K\frac{\log n}{\log\log n}
\]
distinct elements of $S$, for a fixed $K>0$.  Choose a sufficiently small constant $c>0$ and put
\[
N=\left\lfloor\exp(c m\log m)\right\rfloor.
\]
Choose a fixed integer $n_0>\ee^{\ee}$. For $n_0\le n\le N$,
monotonicity gives
\[
K\scale(n)\le K\scale(N)=(Kc+o(1))m.
\]
Taking $c<1/(2K)$ makes this at most $m$ for all sufficiently large $m$.
The finitely many positive $n<n_0$ also have representations of bounded
weight, for example their binary expansions. Increase the threshold for
$m$ so that it exceeds every one of those weights. Thus every integer
$1\le n\le N$ is representable with at most $m$ terms. It follows that
$\ell_m(S)>\lfloor\exp(cm\log m)\rfloor$, and since $\ell_m(S)$ is an
integer, $\ell_m(S)>\exp(cm\log m)$. This proves the lower bound.

\smallskip
\noindent\emph{Upper bound.}
Let
\[
T(N):=|S\cap[1,N]|.
\]
If $b_1^{a_1}\cdots b_r^{a_r}\le N$, then
$a_i\le(\log N)/(\log b_i)$, so
\begin{equation}\label{eq:TN}
T(N)\le C_0(\log N)^r
\end{equation}
for a constant $C_0$ depending only on the bases.  Every representation of an integer at most $N$ as a sum of at most $m$ positive elements of $S$ uses only elements of $S\cap[1,N]$.  The number of multisets of size at most $m$ drawn from $T(N)$ available elements is at most
\[
\binom{T(N)+m}{m}.
\]
Fix $\eta>0$ and set
\[
N=\left\lceil\exp((r-1+\eta/2)m\log m)\right\rceil.
\]
By \eqref{eq:TN} and the standard estimate
$\binom{u+m}{m}\le (\ee(u+m)/m)^m$,
\begin{align*}
\log\binom{T(N)+m}{m}
&\le m\log\!\left(\frac{\ee(T(N)+m)}{m}\right)\\
&\le (r-1)m\log m+r m\log\log m+O(m).
\end{align*}
The right-hand side is strictly smaller than
\[
\log N=(r-1+\eta/2)m\log m+O(1)
\]
for all sufficiently large $m$.  Hence fewer than $N$ sums are available to cover the $N$ integers $1,\ldots,N$, so at least one is missing.  Therefore $\ell_m(S)\le N$. For sufficiently large $m$, the ceiling can be absorbed using the remaining $\eta/2$ margin, giving $N\le\exp((r-1+\eta)m\log m)$. This proves the upper bound in \eqref{eq:ell-bounds}.
\end{proof}

\begin{corollary}\label{cor:ell-21026}
For the Golay monoid $S(2,10,26)$,
\[
\ell_m(S(2,10,26))=\exp(\Theta(m\log m)),
\]
and more specifically
\[
\limsup_{m\to\infty}
\frac{\log \ell_m(S(2,10,26))}{m\log m}\le2.
\]
\end{corollary}

\begin{proof}
The bases $2,10,26$ are multiplicatively independent: from
\[
2^a10^b26^c=2^{a+b+c}5^b13^c=1
\]
with $a,b,c\in\mathbb Z$, unique factorization gives $b=c=0$ and then $a=0$.  Apply Theorem~\ref{thm:ell-general} with $r=3$.
\end{proof}

\begin{remark}
The same theorem applied to $S(2,3)$ gives
\[
\ell_m(S(2,3))=\exp(\Theta(m\log m)).
\]
Thus the very rapid initial values discussed by Livinskyi cannot have asymptotic growth of order $\exp(\Theta(m^2))$.  This is consistent with the modern sparse double-base representation theory that culminated in Theorem~\ref{thm:ksw}.
\end{remark}

\section{Gaps between Hadamard orders}\label{sec:gaps}
Let
\[
n_1<n_2<n_3<\cdots
\]
be the Hadamard orders and, for real $x\ge1$, define, following Brent, Osborn and Smith~\cite{Brent2013},
\[
\gamma(x)=\max\{n_{i+1}-n_i:n_i\le x\}.
\]
A logarithmic estimate $\tau(q)\le\alpha\log_2q+O(1)$ gives a polynomial gap bound by comparing consecutive odd multipliers at a common power of two.  Theorem~\ref{thm:main} gives a subpolynomial version of the same idea.

\begin{theorem}[Uniform quantitative Hadamard gaps]\label{thm:gaps}
There exist $D>0$ and $N\in\N$ such that for every integer $x\ge N$ there
is an integer $d\ge0$ satisfying
\begin{equation}\label{eq:uniform-gap}
d\le\exp(D\scale(x)),\qquad
\forall y\in\N_0,\ y\le x:\quad
\exists h\in\Hcal,\quad y<h\le y+d.
\end{equation}
Consequently the consecutive-order gap function satisfies
\begin{equation}\label{eq:gap}
\gamma(x)\le\exp\!\left(O\!\left(\frac{\log x}{\log\log x}\right)\right)
=x^{o(1)}.
\end{equation}
In particular, for every fixed $\varepsilon>0$, the same uniform assertion
holds with $d\le x^{\varepsilon}$ for all sufficiently large integer $x$.
\end{theorem}

\begin{proof}
If $h\in\Hcal$, Sylvester doubling gives $2h\in\Hcal$. Thus every exponent
above an admissible exponent is also admissible. By Lemma~\ref{lem:existence},
Theorem~\ref{thm:main}, and eventual monotonicity of $\scale$, there is a
nondecreasing real function $F$ on positive integers, with
\[
F(u)=O(\scale(u))\quad(u\to\infty),
\]
such that $2^tq\in\Hcal$ whenever $q$ is positive odd and the integer
$t\ge F(q)$. For example, take the maximum of the finitely many least
exponents with odd multiplier at most $u$; the main bound and a fixed
initial maximum give the stated estimate.

For a large integer $x$, choose $T=\lceil F(2x)\rceil$ and set
$d=2^{T+1}$. Given $0\le y\le x$, let $q$ be the least positive odd
integer with $2^Tq>y$. Consecutive odd integers are two apart, so
\[
q\le y/2^T+2\le x+2\le2x,\qquad 0<2^Tq-y\le2^{T+1}.
\]
Since $T\ge F(q)$, the order $h=2^Tq$ belongs to $\Hcal$. Also
\[
\log d=(T+1)\log2=O(\scale(2x))=O(\scale(x)).
\]
This proves \eqref{eq:uniform-gap}. Taking $y=n_i\le x$ bounds each
consecutive-order gap, giving \eqref{eq:gap} (also for real $x$ by taking
integer parts). Finally $D\scale(x)\le\varepsilon\log x$ eventually,
because $D/\log\log x\to0$. Exponentiation gives $d\le x^{\varepsilon}$.
\end{proof}

\begin{remark}[Relation to Smith's unpublished result]
Brent, Osborn and Smith reported in the 2013 version of their maximal-determinant paper that an unpublished manuscript of W.~D.~Smith proves
\[
\gamma(x)=O(x^{\eps})\qquad\text{for every }\eps>0,
\]
with potentially very large hidden constants~\cite{Brent2013}.  Consequently, the qualitative assertion $\gamma(x)=x^{o(1)}$ should not be presented as historically new.  The significance of Theorem~\ref{thm:gaps}, if the present transfer is not already in the literature, is that it supplies the explicit rate
$\exp(O(\log x/\log\log x))$ from published ingredients.
\end{remark}

\section{A consequence for approximate Hadamard matrices}\label{sec:approx}
Alexeev, Jasper and Mixon~\cite{Alexeev2025} study two relaxations of the Hadamard problem.  Let
\[
u(n)=\min_{M\in O(n)}\|M\|_{\max}
\]
be the smallest possible largest entry of an $n\times n$ orthogonal matrix, and let
\[
\kappa(n)=\min_{A\in\{\pm1\}^{n\times n}}
\frac{\sigma_{\max}(A)}{\sigma_{\min}(A)}
\]
be the best condition number of a sign matrix of order $n$, with value infinity when its smallest singular value is zero. Here $\|M\|_{\max}=\max_{i,j}|M_{ij}|$.
They prove, using the then-current Hadamard-gap exponent $3/23$, that every fixed
$\alpha<17/92$ is admissible in
\[
\kappa(n)\le1+n^{-\alpha}
\]
for all sufficiently large $n$.  They further note that their method would allow every $\alpha<1/4$ under the Hadamard conjecture.  Theorem~\ref{thm:gaps} reaches this limiting exponent without assuming the conjecture.

\begin{corollary}[Flat orthogonal matrices]\label{cor:flat}
For every fixed $0<\eps<1/2$, all sufficiently large $n$ satisfy
\[
u(n)\le \frac{1}{\sqrt n-n^{\eps}}.
\]
\end{corollary}

\begin{proof}
The uniform gap theorem, with parameter $\varepsilon/2$, gives a Hadamard
order $m>n$ with $k=m-n\le n^{\varepsilon/2}<n^{\varepsilon}$ for large
$n$. In particular $k<\sqrt m$.

Here is the submatrix repair estimate of Alexeev--Jasper--Mixon
\cite[Lemmas 3 and 4]{Alexeev2025}. Normalize a Hadamard matrix of order
$m$ by $1/\sqrt m$ and partition the resulting orthogonal matrix as
\[
U=\begin{pmatrix}A&B\\C&D\end{pmatrix},
\]
where $A$ is $k\times k$ and $D$ is $n\times n$. Since
$\|A\|_{\rm op}\le k/\sqrt m<1$, the inverse of $I_k-A$ exists, with norm
at most $(1-k/\sqrt m)^{-1}$. The matrix
\[
M=D+C(I_k-A)^{-1}B
\]
is orthogonal: if $z=(I_k-A)^{-1}Bv$, then
$U(z,v)=(z,Mv)$, and equality of norms implies $\|Mv\|=\|v\|$.
Each row of $C$ and column of $B$ has Euclidean norm $\sqrt{k/m}$, hence
\[
|M_{ij}|\le\frac1{\sqrt m}+
\frac{k/m}{1-k/\sqrt m}=\frac1{\sqrt m-k}
\le\frac1{\sqrt n-n^{\varepsilon}}.
\]
The final denominator is positive for large $n$ because
$0<\varepsilon<1/2$. This proves the claimed bound.
\end{proof}

\begin{corollary}[Approximate Hadamard condition numbers]\label{cor:kappa}
For every $\delta>0$, all sufficiently large $n$ satisfy
\begin{equation}\label{eq:kappa}
\kappa(n)\le1+n^{-1/4+\delta}.
\end{equation}
\end{corollary}

\begin{proof}
The random-rounding and concentration argument in
Alexeev--Jasper--Mixon~\cite[proof of Theorem 1]{Alexeev2025} shows that
the flatness estimate with a fixed $0<\varepsilon<1/2$ gives
\[
\kappa(n)\le1+n^{-\alpha}
\]
eventually for every fixed $0<\alpha<1/4-\varepsilon/2$.
Given any $\delta>0$, put $\delta_0=\min(\delta/2,1/8)$, take
$\varepsilon=\delta_0$, and choose $\alpha=1/4-\delta_0$.
These parameters satisfy the required strict inequalities. Therefore
\[
\kappa(n)\le1+n^{-1/4+\delta_0}\le1+n^{-1/4+\delta}
\]
for sufficiently large $n$, including when $\delta\ge1/4$.
\end{proof}

\begin{remark}
The exponent in Corollary~\ref{cor:kappa} is the limiting value of the Alexeev--Jasper--Mixon gap-based method even if the full Hadamard conjecture is assumed.  Their paper emphasizes that other explicit constructions suggest that the true optimal decay for $\kappa(n)-1$ may be faster; Corollary~\ref{cor:kappa} concerns only what follows from improving the Hadamard-order gaps.
\end{remark}

\section{Effectivity, limitations, and further problems}
The main deduction combines a sparse-representation theorem with a complementary-sequence construction.  That same modularity makes its present limitations clear.

\paragraph{Effective constants.}
Theorem~\ref{thm:ksw} is asymptotic, and Krenn--Suppakitpaisarn--Wagner explicitly note that the finite preprocessing threshold required by their greedy proof can be very large.  Thus Theorem~\ref{thm:main} does not presently provide a competitive numerical construction for a specified moderate $p$.  A useful next result would be an effective version for the particular bases $(2,10)$, with explicit constants and a realistic threshold.  Because the multi-base proof is driven by Diophantine approximation between logarithms of the bases, continued fractions and explicit irrationality measures for $\log 10/\log 2$ are natural tools.

\paragraph{Sharper covering constants.}
For $S(2,10,26)$ we proved only
\[
\exp(c m\log m)\le\ell_m(S)\le\exp((2+o(1))m\log m).
\]
Determining the optimal constants in
\[
\frac{\log\ell_m(S)}{m\log m}
\]
would refine Livinskyi's additive program and might yield useful finite decompositions even without changing the $\Theta(m\log m)$ order.

\paragraph{Beyond a fixed Golay monoid.}
The lower-bound half of the multi-base theorem indicates that $\Theta(\log n/\log\log n)$ is intrinsic to fixed finite collections of multiplicatively independent bases in this representation model.  To beat Theorem~\ref{thm:main} by the same broad strategy one should therefore exploit more than a fixed finite multiplicative monoid: for example, additional primitive Golay lengths, complementary sets of bounded cardinality at many new lengths, or additive constructions that do not reduce to fixed-base power-products.

\paragraph{Reducing the transfer overhead.}
Craigen's theorem converts $r$ complementary pairs into a factor $4^{r+1}$ in the Hadamard order.  Any construction replacing this by $2^{o(r)}$, while retaining sufficiently flexible lengths, would immediately improve the exponent in Theorem~\ref{thm:main}.  Conversely, an obstruction to such a transfer would help identify the real bottleneck.

\section{Formalization and reproducibility}\label{sec:formalization}
The accompanying project, \emph{HadamardBound}, formalizes the main
existence statement, the least exponent, uniform gaps, flat orthogonal
matrices, sign-matrix conditioning, and the stated independent-base
covering bounds in Lean with mathlib:
\begin{center}
\url{https://github.com/L-iet/HadamardBound}.
\end{center}
The mathematical statements in this paper stand independently of the
formalization. The sparse theorem above is presented using its published
source; the Lean project proves the required sparse conclusions by its
own logarithmic-separation, rotation-mesh and greedy arguments from
mathlib. It does not import proof code for the general arbitrary-digit
theorem or its matching infinitely-often lower bound.

All names below have the prefix \path{HadamardMultibase.}. The
unconditional endpoints are exported from
\path{HadamardMultibase/CompleteResults.lean}.
\begin{center}
\small
\begin{tabular}{@{}>{\raggedright\arraybackslash}p{0.35\textwidth}>{\raggedright\arraybackslash}p{0.58\textwidth}@{}}
\toprule
Statement & Lean declaration\\
\midrule
Sparse decomposition & \path{sparseRepresentationBound}\\
Main Hadamard bound & \path{Complete.main_bound}\\
Least admissible exponent & \path{Complete.tau_bound}\\
Uniform quantitative gaps & \path{Complete.quantitative_gap_bound}\\
Subpolynomial gaps & \path{Complete.subpolynomial_gap_bound}\\
Flat orthogonal matrices & \path{Complete.flat_orthogonal_bound}\\
Sign-matrix condition numbers & \path{Complete.optimalConditionNumber_bound}\\
Independent-base covering, lower bound & \path{Complete.powerMonoid_covering_lower}\\
Independent-base covering, upper bound & \path{coveringNumber_upper}\\
Three-base two-sided covering & \texttt{Complete.golayCoveringNumber}\ensuremath{{}_3}\texttt{\_two\_sided}\\
\bottomrule
\end{tabular}
\end{center}
The three-base identifier uses a subscript-three character in the Lean
source. The project uses Lean 4.35.0 and pinned mathlib dependencies.
Its local verification record dated 3 November 2026 reports successful
builds, statement comparison against six independent specifications, and
an independent NanoDa replay of the exported proof. The axiom audit lists
only propositional extensionality, classical choice and quotient soundness.
The six intentional specification placeholders in \path{Challenge.lean}
are excluded from proof imports; \path{Solution.lean} exports completed proofs.

The formal gap predicate uses natural-number starting points exactly as
in \eqref{eq:uniform-gap}. The flatness result requires
$0<\varepsilon<1/2$. The condition-number convention assigns infinity
to singular sign matrices. These domain and convention choices are part
of the assertions, rather than suppressed hypotheses.

Reproduction commands, dependency licences, statement alignment and
recorded hashes are supplied in the repository's \path{README.md},
\path{DEPENDENCIES.md}, \path{paper/ALIGNMENT.md}, and
\path{submission/verification.json}. The project has no additional
proof-code dependency beyond its listed permissively licensed packages.
Formal checking establishes the specified mathematical statements; it
does not establish a historical priority claim.

\begin{thebibliography}{99}
\raggedright

\bibitem{Alexeev2025}
B. Alexeev, J. Jasper, and D. G. Mixon,
\emph{Asymptotically optimal approximate Hadamard matrices},
\href{https://arxiv.org/abs/2511.14653v2}{arXiv:2511.14653v2}, 2025; revised 2026.

\bibitem{Brent2013}
R. P. Brent, J.-A. H. Osborn, and W. D. Smith,
\emph{Lower bounds on maximal determinants of $\{\pm1\}$ matrices via the probabilistic method},
\href{https://arxiv.org/abs/1211.3248v3}{arXiv:1211.3248v3}, 2013.


\bibitem{Craigen1995}
R. Craigen,
\emph{Signed groups, sequences, and the asymptotic existence of Hadamard matrices},
Journal of Combinatorial Theory, Series A \textbf{71} (1995), 241--254.
\href{https://doi.org/10.1016/0097-3165(95)90002-0}{doi:10.1016/0097-3165(95)90002-0}.

\bibitem{DuJiang2024}
C. Du and Y. Jiang,
\emph{Golay complementary sequences of arbitrary length and asymptotic existence of Hadamard matrices},
\href{https://arxiv.org/abs/2401.15381}{arXiv:2401.15381}, 2024.

\bibitem{Golay1961}
M. J. E. Golay,
\emph{Complementary series},
IRE Transactions on Information Theory \textbf{7} (1961), 82--87.
\href{https://doi.org/10.1109/TIT.1961.1057620}{doi:10.1109/TIT.1961.1057620}.

\bibitem{Krenn2020}
D. Krenn, V. Suppakitpaisarn, and S. Wagner,
\emph{On the minimal Hamming weight of a multi-base representation},
Journal of Number Theory \textbf{208} (2020), 168--179.
\href{https://doi.org/10.1016/j.jnt.2019.07.023}{doi:10.1016/j.jnt.2019.07.023}; preprint \href{https://arxiv.org/abs/1808.06330v2}{arXiv:1808.06330v2}.

\bibitem{Livinskyi2012}
I. Livinskyi,
\emph{Asymptotic Existence of Hadamard Matrices},
M.Sc. thesis, University of Manitoba, 2012.

\bibitem{Turyn1974}
R. J. Turyn,
\emph{Hadamard matrices, Baumert--Hall units, four-symbol sequences, pulse compression, and surface wave encodings},
Journal of Combinatorial Theory, Series A \textbf{16} (1974), 313--333.
\href{https://doi.org/10.1016/0097-3165(74)90056-9}{doi:10.1016/0097-3165(74)90056-9}.

\bibitem{Wallis1976}
J. S. Wallis,
\emph{On the existence of Hadamard matrices},
Journal of Combinatorial Theory, Series A \textbf{21} (1976), 188--195.
\href{https://doi.org/10.1016/0097-3165(76)90062-5}{doi:10.1016/0097-3165(76)90062-5}.

\end{thebibliography}

\end{document}
